\subsection{Sub-Phase 5.2: Real-Time Parameter Synchronization and Visualization}
\label{subsec:interactivity}

The definitive feature of the Digital Twin is the bi-directional link between user inputs and the virtual model. As a user modifies a physical or geometric parameter $\vect{\mu}$ (e.g., thickness, material stiffness, or hole radius), the \acrshort{rom} engine performs an instantaneous state reconstruction:
\begin{equation}
    \vect{u}(\vect{\mu}) \approx \mat{V}_r \vect{q}_r(\vect{\mu})
\end{equation}

Because the hyper-reduced solver operates on a generalized coordinate vector $\vect{q}_r$ of minimal dimension (typically $r < 20$), the Newton-Raphson loop converges in under 16ms, enabling a 60FPS interaction rate. The resulting displacement field is then mapped back to the physical NURBS control points, updating the visual mesh in real-time.

\begin{figure}[H]
    \centering
    \includegraphics[width=0.85\textwidth]{digital_twin_light_hero.png}
    \caption{Final Interactive UI: Real-time structural feedback as the user explores the design space via parameter sliders.}
    \label{fig:ui_screenshot}
\end{figure}

The visualization engine includes dynamic features such as:
\begin{itemize}
    \item \textbf{Stress Heatmaps:} Real-time calculation of Von Mises stresses derived from the NURBS derivatives.
    \item \textbf{Deformation Scaling:} Interactive control over the visual amplification of small strains for better qualitative understanding.
    \item \textbf{Comparison Mode:} Overlapping the reduced model result with the baseline linear solution to highlight the impact of geometric nonlinearity.
\end{itemize}

\begin{devnote}
In \texttt{iga-sim-5.2/src/ui/sliders.js}, I implemented a "debounce" logic for the most expensive geometric parameters. This ensures that the ROM engine isn't overwhelmed by redundant calls while the user is rapidly dragging a slider, maintaining a smooth and responsive experience even on mobile devices.
\end{devnote}